\section[Práctica inicial con productos]{Práctica inicial con productos de números naturales}

Antes de ampliar el producto a otros conjuntos numéricos, conviene practicar las ideas más elementales que acabamos de introducir. En estos ejercicios no buscamos todavía trabajar con cálculos complicados, sino reconocer distintas formas de representar un mismo número y distinguir la estructura de las expresiones.

\subsection*{De sumas repetidas a productos}

Escriba cada suma repetida como un producto. No es necesario evaluarlo en un primer paso.

\begin{enumerate}
\item \(3+3+3\)
\item \(5+5+5+5\)
\item \(2+2+2+2+2+2\)
\item \(7+7\)
\item \(10+10+10\)
\item \(1+1+1+1+1\)
\end{enumerate}

Por ejemplo,
\[
4+4+4=3\cdot4.
\]
Observe que el primer factor indica cuántas veces aparece el segundo como sumando.

\subsection*{De productos a sumas repetidas}

Realice ahora el proceso contrario. Escriba cada producto como una suma repetida y luego encuentre un único numeral que represente su valor.

\begin{enumerate}
\item \(2\cdot3\)
\item \(4\cdot2\)
\item \(3\cdot5\)
\item \(5\cdot4\)
\item \(2\cdot10\)
\item \(6\cdot1\)
\end{enumerate}

\subsection*{Evaluar productos}

Encuentre un único numeral que represente el valor de cada expresión.

\begin{enumerate}
\item \(3\cdot4\)
\item \(6\cdot2\)
\item \(5\cdot5\)
\item \(7\cdot3\)
\item \(8\cdot1\)
\item \(9\cdot0\)
\item \(2\cdot4\cdot3\)
\item \(5+2\cdot3\)
\item \(4\cdot2+7\)
\item \(2+3\cdot4+1\)
\end{enumerate}

Recuerde que, cuando aparecen sumas y productos en una misma expresión, el producto se considera agrupado antes que la suma.

\subsection*{Completar expresiones equivalentes}

Complete los espacios de manera que ambos lados de cada igualdad representen el mismo número natural.

\begin{enumerate}
\item \(3\cdot4=\square\)
\item \(5+5+5=\square\cdot5\)
\item \(4\cdot3=3\cdot\square\)
\item \(7\cdot1=\square\)
\item \(9\cdot0=\square\)
\item \(2\cdot(3+4)=2\cdot3+2\cdot\square\)
\end{enumerate}

\subsection*{Reconocer la estructura de una expresión}

En cada caso, indique cuál es la operación principal. Si la expresión principal es una suma, identifique sus sumandos; si contiene productos, identifique también los factores de cada producto.

\begin{enumerate}
\item \(3\cdot5\)
\item \(7+2\cdot4\)
\item \(3\cdot2+5\)
\item \(2+4\cdot3+1\)
\item \(5\cdot(2+3)\)
\item \((2+3)\cdot4\)
\end{enumerate}

El objetivo de estos ejercicios no es únicamente obtener resultados. También debe poder reconocer qué operación está actuando en cada nivel de una expresión y qué papel cumple cada número dentro de ella.

